\chapter{Line Integrals and the Curl}

\epigraph{\textit{To err from the right path is common to Mankind.}}{Sophocles}

%% ================================================================
\section{Work and Line Integrals}
%% ================================================================

We remarked above that the differential form of Gauss' law,
Equations~(II-18) and~(II-23), although it fulfills our goal of relating
a property of the electric field (its divergence) at a point to a known
quantity (the charge density) at the same point, nonetheless falls short
of providing a convenient way to find~$\bfE$.  The reason is that
$\nabla\cdot\bfE = \rho/\epszero$ is (or seems to be) a single
differential equation in three unknowns ($E_x$, $E_y$, $E_z$).  But
there is another feature of electrostatic fields that has not yet played
an explicit role in our discussion and that will yield a relationship
among the components of~$\bfE$.  It will thus provide us with the
crucial last step in obtaining a useful way to calculate fields.  In the
process of examining this question, we shall encounter some of the most
important topics in vector calculus.

The property of electrostatic fields that we shall now begin to
discuss is intimately bound up with the question of work and energy.
You no doubt recall the elementary definition of work as force times
distance.  Thus, in one dimension, if a force $F(x)$ acts from $x=a$
to $x=b$, the work done is, by definition,
\[
  \int_a^b F(x)\dx.
\]

To be able to handle more general situations, we must now introduce
the concept of the \emph{line integral}.

\begin{figure}[ht]
\centering
\includegraphics[width=0.6\textwidth]{figures/ch03/fig_III_1.png}
\caption{}
\label{fig:III-1}
\end{figure}

Suppose we have a curve~$C$ in three dimensions (Figure~III-1) and
suppose the curve is \emph{directed}.  By this we mean that we put an
arrow on the curve and say ``This is the positive direction.''  Let~$s$
be the arc length measured along the curve from some arbitrary point on
it with $s=s_1$ at a point~$P_1$ and $s=s_2$ at~$P_2$.  Suppose
further that we have a function $f(x,y,z)$ defined everywhere on~$C$.
Now let us subdivide the portion of~$C$ between $P_1$ and $P_2$
arbitrarily into $N$~sections.  Figure~III-1 shows an example of such a
subdivision for $N=4$.  Next, join successive subdivision points by
chords, a typical one of which, say the $l$th, has length~$\Delta s_l$.
Now evaluate $f(x,y,z)$ at $(x_l,y_l,z_l)$, which is any point on the
$l$th subdivision of the curve, and form the product
$f(x_l,y_l,z_l)\,\Delta s_l$.  Doing this for each of the $N$~segments
of~$C$, we form the sum
\[
  \sum_{l=1}^{N} f(x_l,y_l,z_l)\,\Delta s_l.
\]

By definition, the line integral of $f(x,y,z)$ along the curve~$C$ is
the limit of this sum as the number of subdivisions $N$ approaches
infinity and the length of each chord approaches zero:
\begin{equation}
  \int_C f(x,y,z)\ds = \lim_{\substack{N\to\infty\\ \text{each }\Delta s_l\to 0}}
  \sum_{l=1}^{N} f(x_l,y_l,z_l)\,\Delta s_l.
\end{equation}

To evaluate the line integral, we need to know the path~$C$.
Usually the most convenient way to specify this path is parametrically
in terms of the arc length parameter~$s$.  Thus, we write $x=x(s)$,
$y=y(s)$, and $z=z(s)$.  In such a situation the line integral can be
reduced to an ordinary definite integral:
\begin{equation}
  \int_C f(x,y,z)\ds = \int_{s_1}^{s_2} f[x(s),\,y(s),\,z(s)]\ds.
\end{equation}

An example of a line integral will be helpful here.  For simplicity
let us work in two dimensions and evaluate
\[
  \int_C (x+y)\ds,
\]
where $C$~is the straight line from the origin to the point whose
coordinates are~$(1,1)$ (Figure~III-2).  If $(x,y)$ are the coordinates
of

\begin{figure}[ht]
\centering
\includegraphics[width=0.5\textwidth]{figures/ch03/fig_III_2.png}
\caption{}
\label{fig:III-2}
\end{figure}

\noindent any point~$P$ on~$C$ and if $s$~is the arc length measured
from the origin, then $x = s/\sqrt{2}$ and $y = s/\sqrt{2}$.  Hence,
$x+y = 2s/\sqrt{2} = \sqrt{2}\,s$.  Thus,
\begin{equation}
  \int_C (x+y)\ds = \sqrt{2}\int_0^{\sqrt{2}} s\ds = \sqrt{2}.
\end{equation}

\begin{figure}[ht]
\centering
\includegraphics[width=0.5\textwidth]{figures/ch03/fig_III_3.png}
\caption{}
\label{fig:III-3}
\end{figure}

Let us integrate this same function $(x+y)$ from $(0,0)$ to $(1,1)$
along another path as shown in Figure~III-3.  Here we break the
integration into two parts, one along~$C_1$ and the second along~$C_2$.
On $C_1$ we have $x=s$ and $y=0$.  Thus, on $C_1$, $x+y=s$, and so
\[
  \int_{C_1}(x+y)\ds = \int_0^1 s\ds = \tfrac{1}{2}.
\]

Along $C_2$, $x=1$ and $y=s$ [note that the arc length on this
segment of the path is measured from the point~$(1,0)$].  It follows
then that
\[
  \int_{C_2}(x+y)\ds = \int_0^1 (1+s)\ds = \tfrac{3}{2}.
\]

Adding the results for the two segments, we find
\[
  \int_C (x+y)\ds = \int_{C_1}(x+y)\ds + \int_{C_2}(x+y)\ds
  = \tfrac{1}{2}+\tfrac{3}{2} = 2.
\]

The lesson to be learned is this: the value of a line integral can
(indeed, usually does) depend on the path of integration.


%% ================================================================
\section{Line Integrals Involving Vector Functions}
%% ================================================================

Although the preceding discussion tells us what a line integral is,
the kind of line integral we must deal with here has a feature not yet
mentioned.  You will recall that we introduced our discussion of line
integrals with the concept of work.  Work, in the most elementary
sense, is force times displacement.  That this needs elaboration becomes
clear when we recognize that both force and displacement are vectors.

Thus, consider some path~$C$ in three dimensions (Figure~III-4).
Let us suppose that under the action of a force an object moves on this
path from $s_1$ to~$s_2$.  At any point~$P$ on the curve let the force

\begin{figure}[ht]
\centering
\includegraphics[width=0.6\textwidth]{figures/ch03/fig_III_4.png}
\caption{}
\label{fig:III-4}
\end{figure}

\noindent acting be designated $\mathbf{f}(x,y,z)$.  The component of
$\mathbf{f}$ that does work is, by definition, only that one which acts
\emph{along} the curve, that is, the tangential component.  Let
$\hat{\mathbf{t}}$ denote a unit vector that is tangent to the curve
at~$P$.\footnote{$\hat{\mathbf{t}}$ is a function of $x$, $y$, and $z$
and should really be written $\hat{\mathbf{t}}(x,y,z)$.  We write
simply $\hat{\mathbf{t}}$ to avoid complicating the notation.}  Then
the work done by the force in moving the object from $s_1$ to~$s_2$
along the curve~$C$ is
\begin{equation}
  W = \int_C \mathbf{f}(x,y,z)\cdot\hat{\mathbf{t}}\ds,
\end{equation}
where it is understood, of course, that the integration begins at
$s=s_1$ and ends at $s=s_2$.  The new feature of this integral is that
the integrand is the dot product of two vector functions.  To be able to
handle such a line integral, we must know how to find
$\hat{\mathbf{t}}$, and it is to this problem that we now turn.

Consider an arbitrary curve~$C$ (Figure~III-5) parametrized by its
arc length.  At some point~$s$ on the curve we have $x=x(s)$, $y=y(s)$,
and $z=z(s)$.  At another point $s+\Delta s$ we have
$x+\Delta x = x(s+\Delta s)$, $y+\Delta y = y(s+\Delta s)$, and
$z+\Delta z = z(s+\Delta s)$.  Thus,

\begin{figure}[ht]
\centering
\includegraphics[width=0.6\textwidth]{figures/ch03/fig_III_5.png}
\caption{}
\label{fig:III-5}
\end{figure}

\noindent the chord joining the two points on the curve directed from
the first to the second is the vector
$\Delta\bfr \equiv \bfi\,\Delta x + \bfj\,\Delta y + \bfk\,\Delta z$,
where
\begin{align*}
  \Delta x &= x(s+\Delta s) - x(s), \\
  \Delta y &= y(s+\Delta s) - y(s), \\
  \Delta z &= z(s+\Delta s) - z(s).
\end{align*}

If we now divide this vector by~$\Delta s$, we get
\[
  \frac{\Delta\bfr}{\Delta s}
  = \bfi\,\frac{\Delta x}{\Delta s}
  + \bfj\,\frac{\Delta y}{\Delta s}
  + \bfk\,\frac{\Delta z}{\Delta s}.
\]

Taking the limit of this as $\Delta s$ approaches zero yields
\[
  \bfi\,\frac{dx}{ds} + \bfj\,\frac{dy}{ds} + \bfk\,\frac{dz}{ds},
\]
and we assert that this is~$\hat{\mathbf{t}}$.  To begin with, it's
clear that as $\Delta s\to 0$, the vector $\Delta\bfr$ becomes tangent
to the curve at~$s$.  Further, in the limit $\Delta s\to 0$, we see
that $|\Delta\bfr|\to\Delta s$.  Hence, in the limit the magnitude of
this quantity is~1.  It follows then that we can make the identification
\begin{equation}
  \hat{\mathbf{t}}(s)
  = \bfi\,\frac{dx}{ds} + \bfj\,\frac{dy}{ds} + \bfk\,\frac{dz}{ds}.
\end{equation}

If we return now to the expression for work~$W$ and use this formula
for~$\hat{\mathbf{t}}$, we find
\begin{equation}
  W = \int_C \mathbf{f}(x,y,z)\cdot
      \left[\bfi\,\frac{dx}{ds}+\bfj\,\frac{dy}{ds}
            +\bfk\,\frac{dz}{ds}\right]ds
    = \int_C (f_x\dx + f_y\dy + f_z\dz).
\end{equation}

This is a formal expression; often, to carry out the integration, it is
useful to restore the $ds$ as the following example illustrates.
Consider
\[
  \mathbf{f}(x,y,z) = \bfi\,y - \bfj\,x
\]
and the path shown in Figure~III-6(a).  To calculate
$\int_C(\mathbf{f}\cdot\hat{\mathbf{t}})\ds$ in

\begin{figure}[ht]
\centering
\includegraphics[width=0.5\textwidth]{figures/ch03/fig_III_6a.png}
\caption{}
\label{fig:III-6a}
\end{figure}

\noindent this case, we break the path~$C$ into three parts, $C_1$,
$C_2$, and $C_3$ as shown.  Since $f_z=0$, we have
\begin{align*}
  \int_C \mathbf{f}\cdot\hat{\mathbf{t}}\ds
    &= \int_C f_x\dx + f_y\dy \\
    &= \int_C y\dx - x\dy.
\end{align*}

Now, on~$C_1$, $y=0$ and $dy=0$, so there is no contribution to the
integral.  Similarly, on $C_3$ we have $x=0$ and $dx=0$, and again the
result is zero.  Thus, the only contribution to the integral over~$C$
can come from~$C_2$.  Restoring the $ds$, we have
\[
  \int_C \left(y\,\frac{dx}{ds} - x\,\frac{dy}{ds}\right)ds.
\]

\begin{figure}[ht]
\centering
\includegraphics[width=0.5\textwidth]{figures/ch03/fig_III_6b.png}
\caption{}
\label{fig:III-6b}
\end{figure}

But $(1-x)/s = \cos 45^\circ = 1/\sqrt{2}$ and $y/s = \sin 45^\circ
= 1/\sqrt{2}$ [Figure~III-6(b)].  Thus,
\[
  x = 1 - \frac{s}{\sqrt{2}} \;\Rightarrow\; \frac{dx}{ds}
    = -\frac{1}{\sqrt{2}} \qquad
  y = \frac{s}{\sqrt{2}} \;\Rightarrow\; \frac{dy}{ds}
    = \frac{1}{\sqrt{2}}
  \qquad 0\le s\le\sqrt{2}.
\]

Hence, the integral is
\begin{align*}
  \int_0^{\sqrt{2}} &\left[\frac{s}{\sqrt{2}}
    \left(-\frac{1}{\sqrt{2}}\right)
    - \left(1-\frac{s}{\sqrt{2}}\right)\frac{1}{\sqrt{2}}\right]ds \\
  &= -\frac{1}{\sqrt{2}}\int_0^{\sqrt{2}}ds = -1.
\end{align*}

As a second example of a line integral involving a vector function, let
\[
  \mathbf{f}(x,y,z) = \bfi\,x^2 - \bfj\,xy,
\]
and take~$C$ to be the quarter circle of radius~$R$ oriented as shown
in Figure~III-6(c).  We then have
\[
  \int_C \mathbf{f}\cdot\hat{\mathbf{t}}\ds
    = \int_C x^2\dx - xy\dy.
\]

\begin{figure}[ht]
\centering
\includegraphics[width=0.45\textwidth]{figures/ch03/fig_III_6c.png}
\caption{}
\label{fig:III-6c}
\end{figure}

Letting $x=R\cos\theta$, $y=R\sin\theta$, we find this integral becomes
\begin{align*}
  \int_0^{\pi/2} [R^2\cos^2\theta(-R\sin\theta)
    - R^2\sin\theta\cos\theta(R\cos\theta)]\,d\theta
  &= -2R^3\int_0^{\pi/2}\cos^2\theta\,\sin\theta\,d\theta \\
  &= -2R^3/3.
\end{align*}


%% ================================================================
\section{Path Independence}
%% ================================================================

In a line integral the path of integration is one of the ingredients
which determines the very function we integrate.  It isn't remarkable,
then, that the value of the integral can depend on the path of
integration.  What is remarkable is that, under some conditions, the
value of the integral does \emph{not} depend on the path!

We show how this path independence comes about in the case of the
Coulomb force.  Let a charge~$q_0$ be fixed at the origin and let
another charge~$q$ be situated at $(x,y,z)$ (Figure~III-7).  The
Coulomb force on~$q$ is
\begin{equation}
  \bfF = \frac{1}{4\pi\epszero}\,\frac{qq_0}{r^2}\,\hat{\mathbf{u}},
  \label{eq:III-1}
\end{equation}
where $r=(x^2+y^2+z^2)^{1/2}$ is the distance between the two charges
and $\hat{\mathbf{u}}$~is a unit vector pointing from $q_0$ to~$q$.
With this arrangement $\hat{\mathbf{u}}$~is clearly in the radial
direction.  Even more

\begin{figure}[ht]
\centering
\includegraphics[width=0.6\textwidth]{figures/ch03/fig_III_7.png}
\caption{}
\label{fig:III-7}
\end{figure}

\noindent clearly, the radial vector~$\bfr$ is in the radial direction.
Thus, we have $\hat{\mathbf{u}}=\bfr/r=(\bfi\,x+\bfj\,y+\bfk\,z)/r$,
and so
\[
  \bfF = \frac{qq_0}{4\pi\epszero}\,
         \frac{\bfi\,x+\bfj\,y+\bfk\,z}{r^3}.
\]

Thus,
\[
  \bfF\cdot\hat{\mathbf{t}}\ds
  = F_x\dx + F_y\dy + F_z\dz
  = \frac{qq_0}{4\pi\epszero}\,\frac{x\dx + y\dy + z\dz}{r^3}.
\]

The trick now is to use the relationship
\[
  r^2 = x^2 + y^2 + z^2.
\]

Taking differentials in this equation and dividing by a factor of~2
yields
\[
  x\dx + y\dy + z\dz = r\,dr,
\]
so that
\[
  \bfF\cdot\hat{\mathbf{t}}\ds
  = \frac{qq_0}{4\pi\epszero}\,\frac{r\,dr}{r^3}
  = \frac{qq_0}{4\pi\epszero}\,\frac{dr}{r^2}.
\]

Suppose now that the charge~$q$ moves from a point~$P_1$ at a
distance~$r_1$ from the origin to a point~$P_2$ at a distance~$r_2$,
over some path~$C$ connecting the two points (Figure~III-8).  Then
\begin{equation}
  \int_C \bfF\cdot\hat{\mathbf{t}}\ds
  = \frac{qq_0}{4\pi\epszero}\int_{r_1}^{r_2}\frac{dr}{r^2}
  = \frac{qq_0}{4\pi\epszero}
    \left(\frac{1}{r_1}-\frac{1}{r_2}\right).
\end{equation}

\begin{figure}[ht]
\centering
\includegraphics[width=0.6\textwidth]{figures/ch03/fig_III_8.png}
\caption{}
\label{fig:III-8}
\end{figure}

Notice that to get this result, we haven't had to specify $C$ in any
way whatever; we'd get the same answer for \emph{any} path connecting
$P_1$ and~$P_2$.  This, of course, proves that the line integral
\[
  \int_C \bfF\cdot\hat{\mathbf{t}}\ds
\]
with $\bfF$ given by Equation~(\ref{eq:III-1}) is path-independent,
but the result, so far, has been established only for the Coulomb force
on~$q$ due to a single charge~$q_0$ [Equation~(\ref{eq:III-1})].  If
there are many charges $q_1,q_2,\dots,q_N$, then the total force on~$q$
is $\bfF_1+\bfF_2+\cdots+\bfF_N$, where $\bfF_l$~is the Coulomb force
on~$q$ due to the $l$th charge~$q_l$.  Hence,
\[
  \int_C \bfF\cdot\hat{\mathbf{t}}\ds
  = \int_C \bfF_1\cdot\hat{\mathbf{t}}\ds
  + \cdots
  + \int_C \bfF_N\cdot\hat{\mathbf{t}}\ds.
\]

Now the discussion given above shows that each term of this sum is path
independent; hence, so is the sum itself.  (All this, of course, is
merely an application of the superposition principle.)  To phrase this
result in terms of the field requires one last trivial step: Since
$\bfF=q\bfE$, it follows that
$q\int_C\bfE\cdot\hat{\mathbf{t}}\ds$ is path independent, whence
$\int_C\bfE\cdot\hat{\mathbf{t}}\ds$ is also.  Strange to say, it is
this fact that will enable us eventually to convert
$\nabla\cdot\bfE=\rho/\epszero$ into a more useful equation.

If you examine the foregoing discussion carefully, you'll see that
the fact that the Coulomb force varies inversely as the square of~$r$
has nothing whatever to do with the path independence of the line
integral.  The path independence rests solely on two properties of the
Coulomb force: (1) It depends only on the distance between the two
particles, and (2) it acts along the line joining them.

Any force~$\bfF$ with these two properties is called a
\emph{central force}, and $\int_C\bfF\cdot\hat{\mathbf{t}}\ds$ is
independent of path for \emph{any} central
force.\footnote{Our having illustrated path independence with a central
force may give the erroneous impression that \emph{only} central forces
have path-independent line integrals.  That is certainly \emph{not}
true; many functions which are not central forces have path-independent
line integrals.  Later we'll develop a simple criterion for identifying
such functions.}

One further step pertaining to path independence can be taken here.
If
\[
  \int_C \bfF\cdot\hat{\mathbf{t}}\ds
\]
is independent of path, then
\[
  \int_{C_1}\bfF\cdot\hat{\mathbf{t}}\ds
  = \int_{C_2}\bfF\cdot\hat{\mathbf{t}}\ds,
\]
where, as indicated in Figure~III-9, $C_1$ and $C_2$ are two different
arbitrary paths connecting the two points $P_1$ and $P_2$ and directed
as shown in the figure.  Now if instead of integrating along~$C_1$ from
$P_1$ to~$P_2$, we go the other way, we simply change the sign of the
line integral; that is,
\[
  \int_{-C_1}\bfF\cdot\hat{\mathbf{t}}\ds
  = -\int_{C_1}\bfF\cdot\hat{\mathbf{t}}\ds,
\]

\begin{figure}[ht]
\centering
\includegraphics[width=0.5\textwidth]{figures/ch03/fig_III_9.png}
\caption{}
\label{fig:III-9}
\end{figure}

\noindent where $-C_1$ merely indicates that the integration is to be
carried out along~$C_1$ from $P_2$ to~$P_1$.  Thus,
\[
  \int_{C_2}\bfF\cdot\hat{\mathbf{t}}\ds
  = -\int_{-C_1}\bfF\cdot\hat{\mathbf{t}}\ds
\]
or
\[
  \int_{-C_1+C_2}\bfF\cdot\hat{\mathbf{t}}\ds = 0.
\]

But $-C_1+C_2$ is just the closed loop from $P_1$ to $P_2$ and back,
as shown in Figure~III-10.  Thus, if
$\int\bfF\cdot\hat{\mathbf{t}}\ds$ is independent of path, then
\[
  \oint \bfF\cdot\hat{\mathbf{t}}\ds = 0,
\]

\begin{figure}[ht]
\centering
\includegraphics[width=0.5\textwidth]{figures/ch03/fig_III_10.png}
\caption{}
\label{fig:III-10}
\end{figure}

\noindent where $\oint$ is the standard notation for a line integral
around a closed path.  It follows that if $\bfE$~is an electrostatic
field, we can write
\begin{equation}
  \oint \bfE\cdot\hat{\mathbf{t}}\ds = 0.
  \label{eq:III-2}
\end{equation}

The term ``circulation'' is often given to the path integral around a
closed curve of the tangential component of a vector function.  Thus we
have demonstrated that the circulation of the electrostatic field is
zero.  In what follows we'll call this the circulation law.


%% ================================================================
\section{The Curl}
%% ================================================================

If we are given some vector function $\bfF(x,y,z)$ and asked, ``Could
this be an electrostatic field?'' we can, in principle, provide an
answer.  If
\[
  \oint \bfF\cdot\hat{\mathbf{t}}\ds \ne 0
\]
over even \emph{one} path, then $\bfF$ cannot be an electrostatic
field.  If
\[
  \oint \bfF\cdot\hat{\mathbf{t}}\ds = 0
\]
over \emph{every} closed path, then $\bfF$ can (but does not have to)
be an electrostatic field.

Clearly, this criterion is not easy to apply since we must be sure the
circulation of~$\bfF$ is zero over all possible paths.  To develop a
more useful criterion, we proceed much as we did in dealing with Gauss'
law, which, like the circulation law, is an expression involving an
integral over the electric field.  Gauss' law is more useful in the
differential form [Equations~(II-18) and~(II-23)] obtained by
considering the ratio of flux to volume for ever-decreasing surfaces.
We now treat the circulation law in the same spirit and attempt to find
the differential form of Equation~(\ref{eq:III-2}).  To stress the
generality of our analysis and results, we deal with an arbitrary
function $\bfF(x,y,z)$ and specialize to $\bfE(x,y,z)$ at a later
stage in the development.

Let us consider the circulation of~$\bfF$ over a small rectangle
parallel to the $xy$-plane, with sides $\Delta x$ and $\Delta y$ and
with the point $(x,y,z)$ at the center [Figure~III-11(a)].  As shown in
Figure~III-11(b), we carry out the path integration in a
counterclockwise direction looking down at the $xy$-plane.  The line
integral is broken up into four parts: $C_B$~(bottom), $C_R$~(right),
$C_T$~(top), and $C_L$~(left).  Since the rectangle is small
(eventually we shall take the limit as it shrinks down to zero), we'll
approximate the integral over each segment by
$\bfF\cdot\hat{\mathbf{t}}$ evaluated at the center of the segment,
multiplied by the length of the
segment.\footnote{Reread footnote~9 of Chapter~II and then give an
argument in support of this approximation.}

\begin{figure}[ht]
\centering
\includegraphics[width=0.6\textwidth]{figures/ch03/fig_III_11a.png}
\caption{}
\label{fig:III-11a}
\end{figure}

\begin{figure}[ht]
\centering
\includegraphics[width=0.5\textwidth]{figures/ch03/fig_III_11b.png}
\caption{}
\label{fig:III-11b}
\end{figure}

Taking $C_B$ first, we have
\begin{equation}
  \int_{C_B}\bfF\cdot\hat{\mathbf{t}}\ds
  = \int_{C_B}F_x\dx
  \approx F_x\!\left(x,\,y-\frac{\Delta y}{2},\,z\right)\Delta x.
  \label{eq:III-3a}
\end{equation}

Over $C_T$ we find
\begin{equation}
  \int_{C_T}\bfF\cdot\hat{\mathbf{t}}\ds
  = \int_{C_T}F_x\dx
  \approx -F_x\!\left(x,\,y+\frac{\Delta y}{2},\,z\right)\Delta x.
  \label{eq:III-3b}
\end{equation}

The negative sign here is required by the fact that
\[
  \int_{C_T}F_x\dx = \int_{C_T}F_x\,\frac{dx}{ds}\ds
\]
and $dx/ds=-1$ over~$C_T$.  Adding Equation~(\ref{eq:III-3a}) and
(\ref{eq:III-3b}), we find
\begin{align*}
  \int_{C_T+C_B}(\bfF\cdot\hat{\mathbf{t}})\ds
  &\approx -\left[F_x\!\left(x,\,y+\frac{\Delta y}{2},\,z\right)
           - F_x\!\left(x,\,y-\frac{\Delta y}{2},\,z\right)\right]\Delta x.
\end{align*}

The factor $\Delta x\,\Delta y$ is clearly the area $\Delta S$ of the
rectangle.  Thus,
\begin{equation}
  \frac{1}{\Delta S}\int_{C_T+C_B}(\bfF\cdot\hat{\mathbf{t}})\ds
  \approx -\frac{F_x\!\left(x,\,y+\frac{\Delta y}{2},\,z\right)
           - F_x\!\left(x,\,y-\frac{\Delta y}{2},\,z\right)}{\Delta y}.
  \label{eq:III-4}
\end{equation}

Exactly the same sort of analysis applied to the left and right sides
of the rectangle ($C_L$ and $C_R$) results in
\begin{equation}
  \frac{1}{\Delta S}\int_{C_L+C_R}(\bfF\cdot\hat{\mathbf{t}})\ds
  \approx \frac{F_y\!\left(x+\frac{\Delta x}{2},\,y,\,z\right)
           - F_y\!\left(x-\frac{\Delta x}{2},\,y,\,z\right)}{\Delta x}.
  \label{eq:III-5}
\end{equation}

Adding Equations~(\ref{eq:III-4}) and~(\ref{eq:III-5}) and taking the
limit as $\Delta S$ shrinks down about a point $(x,y,z)$ (in which case
$\Delta x$ and $\Delta y\to 0$ as well), we get
\begin{equation}
  \lim_{\substack{\Delta S\to 0 \\ \text{about }(x,y,z)}}
  \frac{1}{\Delta S}\oint \bfF\cdot\hat{\mathbf{t}}\ds
  = \pd{F_y}{x} - \pd{F_x}{y},
  \label{eq:III-6}
\end{equation}
where $\oint$ is our semicomical notation meaning the circulation
around the little rectangle.

You may wonder about the generality and uniqueness of this result
since it is obtained using a path of integration that is special in two
ways: first, it is a rectangle, and second, it is parallel to the
$xy$-plane.  If the path were not a rectangle, but a plane curve of
arbitrary shape, it would not affect our result (see Problems~III-2
and~III-30).  But our result definitely \emph{does} depend on the
special orientation of the path of integration.  The choice of
orientation made above clearly suggests two others, and they are shown
in Figure~III-12(a) and~(b) along with the result of calculating
\[
  \lim_{\substack{\Delta S\to 0 \\ \text{about }(x,y,z)}}
  \frac{1}{\Delta S}\oint \bfF\cdot\hat{\mathbf{t}}\ds
\]
for each.

\begin{figure}[ht]
\centering
\includegraphics[width=0.6\textwidth]{figures/ch03/fig_III_12a.png}
\caption{}
\label{fig:III-12a}
\end{figure}

\begin{figure}[ht]
\centering
\includegraphics[width=0.6\textwidth]{figures/ch03/fig_III_12b.png}
\caption{}
\label{fig:III-12b}
\end{figure}

Each of these three paths is named in honor of the vector normal to
the enclosed area.  The convention we use is this: Trace the curve~$C$
so that the enclosed area is always to the left [Figure~III-13(a)].
Then choose the normal so that it points ``up'' in the di-

\begin{figure}[ht]
\centering
\includegraphics[width=0.5\textwidth]{figures/ch03/fig_III_13a.png}
\caption{}
\label{fig:III-13a}
\end{figure}

\begin{figure}[ht]
\centering
\includegraphics[width=0.5\textwidth]{figures/ch03/fig_III_13b.png}
\caption{}
\label{fig:III-13b}
\end{figure}

\noindent rection shown in the picture.  This convention is sometimes
called the right-hand rule, for if the right hand is oriented so that
the fingers curl in the direction in which the curve is traced, the
thumb, extended, points in the direction of the normal
[Figure~III-13(b)].  Using the right-hand rule, we have the following:

\medskip
Calculating $\displaystyle\lim_{\Delta S\to 0}\oint
\bfF\cdot\hat{\mathbf{t}}\ds/\Delta S$

\begin{equation}
\left.
\begin{aligned}
  &\text{for a path whose normal is } \bfi,
    &\quad &\pd{F_z}{y} - \pd{F_y}{z}, \\[4pt]
  &\text{for a path whose normal is } \bfj,
    &\quad &\pd{F_x}{z} - \pd{F_z}{x}, \\[4pt]
  &\text{for a path whose normal is } \bfk,
    &\quad &\pd{F_y}{x} - \pd{F_x}{y}.
\end{aligned}
\;\right\}
\label{eq:III-7a}
\end{equation}

It turns out that these three quantities are the Cartesian components
of a vector.  To this vector we give the name ``curl of~$\bfF$,'' which
we write curl~$\bfF$.  Thus, we have
\begin{equation}
  \text{curl}\;\bfF
  = \bfi\left(\pd{F_z}{y}-\pd{F_y}{z}\right)
  + \bfj\left(\pd{F_x}{z}-\pd{F_z}{x}\right)
  + \bfk\left(\pd{F_y}{x}-\pd{F_x}{y}\right).
  \label{eq:III-7b}
\end{equation}

This expression is often (indeed, usually) given as the definition of
the curl, but we prefer to regard it as merely the form of the curl in
Cartesian coordinates.  We shall define the curl as the limit

\begin{figure}[ht]
\centering
\includegraphics[width=0.5\textwidth]{figures/ch03/fig_III_14.png}
\caption{}
\label{fig:III-14}
\end{figure}

\noindent of circulation to area as the area tends to zero.  To be
precise, let $\oint_{C_n}\bfF\cdot\hat{\mathbf{t}}\ds$ be the
circulation of~$\bfF$ about some path whose normal is~$\nhat$ as shown
in Figure~III-14.  Then by definition
\begin{equation}
  \nhat\cdot\text{curl}\;\bfF
  = \lim_{\substack{\Delta S\to 0 \\ \text{about }(x,y,z)}}
    \frac{1}{\Delta S}\oint_{C_n}\bfF\cdot\hat{\mathbf{t}}\ds.
  \label{eq:III-8}
\end{equation}

By taking $\nhat$ successively equal to $\bfi$, $\bfj$, and $\bfk$, we
get back the results given in Equation~(\ref{eq:III-7b}).  Since this
limit will, in general, have different values for different points
$(x,y,z)$, the curl of~$\bfF$ is a \emph{vector function} of
position.\footnote{The word \emph{rotation} (abbreviated ``rot,''
amusingly enough) was once used for what we now call the \emph{curl}.
Though the term has long since dropped out of use, a related one
survives: If $\text{curl}\;\bfF=0$, the function~$\bfF$ is said to be
\emph{irrotational}.}  Note incidentally that although in our work we
always assumed that the area enclosed by the path of integration was a
plane, this need not be the case.  Since the curl is defined in terms of
a limit in which the enclosed surface shrinks to zero about some point,
in the final stages of this limiting process the enclosed surface is
infinitesimally close to a plane, and all our considerations apply.

Since it is undoubtedly beyond the powers of a mere mortal to
remember the expression given above for curl~$\bfF$ in Cartesian
coordinates [Equation~(\ref{eq:III-7b})], it is fortunate that there is
a mnemonic device to fall back on.  If the three-by-three
determinant\footnote{A mathematician would object to this since,
strictly speaking, a determinant cannot contain either vectors or
operators.  We aren't doing any serious damage, however, because our
``determinant'' is merely a memory aid.}
\[
  \begin{vmatrix}
    \bfi & \bfj & \bfk \\[3pt]
    \partial/\partial x & \partial/\partial y & \partial/\partial z \\[3pt]
    F_x & F_y & F_z
  \end{vmatrix}
\]
is expanded (most conveniently in minors of the first row) and if
certain ``products'' are interpreted as partial derivatives [for example,
$(\partial/\partial x)F_y = \partial F_y/\partial x$], the result will
be identical with the one given in Equation~(\ref{eq:III-7b}).  Thus,
the anguish of remembering the form of curl~$\bfF$ in Cartesian
coordinates can be replaced by the pain of remembering how to expand a
three-by-three determinant.  \emph{Chacun \`a son go\^ut.}

As an example of calculating the curl, consider the vector function
\[
  \bfF(x,y,z) = \bfi\,xz + \bfj\,yz - \bfk\,y^2.
\]
We have
\begin{align*}
  \text{curl}\;\bfF
  &= \begin{vmatrix}
       \bfi & \bfj & \bfk \\
       \partial/\partial x & \partial/\partial y & \partial/\partial z \\
       xz & yz & -y^2
     \end{vmatrix} \\
  &= \bfi(-2y-y) + \bfj(x-0) + \bfk(0-0) \\
  &= -3\bfi\,y + \bfj\,x.
\end{align*}

You may have noticed that the curl operator can be written in terms
of the del notation we introduced earlier.  You can verify for yourself
that
\[
  \text{curl}\;\bfF = \nabla\times\bfF,
\]
which is read ``del cross~$\bfF$.''  Henceforth, we shall always use
$\nabla\times\bfF$ to indicate the curl.


%% ================================================================
\section{The Curl in Cylindrical and Spherical Coordinates}
%% ================================================================

To obtain the form of $\nabla\times\bfF$ in other coordinate systems,
we proceed as we did above in finding the Cartesian form, merely
modifying the paths of integration appropriately.  As an example, using
the path shown in Figure~III-15(a) will yield the $z$-component of
$\nabla\times\bfF$ in cylindrical
coordinates.\footnote{In deriving the Cartesian form of
$\nabla\times\bfF$, each segment of each path of integration (see
Figures~III-11 and~III-12) was of the form $x=\text{constant}$,
$y=\text{constant}$, or $z=\text{constant}$.  Similarly, in deriving the
cylindrical form, each segment of each path is of the form
$r=\text{constant}$, $\theta=\text{constant}$, or $z=\text{constant}$.}
Note that we trace the curve in

\begin{figure}[ht]
\centering
\includegraphics[width=0.6\textwidth]{figures/ch03/fig_III_15a.png}
\caption{}
\label{fig:III-15a}
\end{figure}

\begin{figure}[ht]
\centering
\includegraphics[width=0.5\textwidth]{figures/ch03/fig_III_15b.png}
\caption{}
\label{fig:III-15b}
\end{figure}

\noindent accordance with the right-hand rule given in the previous
section.  Viewing the path from above [as we do in Figure~III-15(b)],
the line integral of
$\bfF(r,\theta,z)\cdot\hat{\mathbf{t}}$ along the segment of path
marked~1 is
\[
  \int_{C_1}\bfF\cdot\hat{\mathbf{t}}\ds
  \approx F_r\!\left(r,\,\theta-\frac{\Delta\theta}{2},\,z\right)
    \Delta r,
\]
while along segment~3 it is
\[
  \int_{C_3}\bfF\cdot\hat{\mathbf{t}}\ds
  \approx -F_r\!\left(r,\,\theta+\frac{\Delta\theta}{2},\,z\right)
    \Delta r.
\]

The area enclosed by the path is $r\,\Delta r\,\Delta\theta$, so
\[
  \frac{1}{\Delta S}\int_{C_1+C_3}\bfF\cdot\hat{\mathbf{t}}\ds
  \approx -\frac{\Delta r}{r\,\Delta r\,\Delta\theta}
    \left[F_r\!\left(r,\,\theta+\frac{\Delta\theta}{2},\,z\right)
    - F_r\!\left(r,\,\theta-\frac{\Delta\theta}{2},\,z\right)\right].
\]

In the limit as $\Delta r$ and $\Delta\theta$ tend to zero, this
becomes
\[
  -\frac{1}{r}\pd{F_r}{\theta}
\]
evaluated at the point $(r,\theta,z)$.

Along segment~2 we find
\[
  \int_{C_2}\bfF\cdot\hat{\mathbf{t}}\ds
  \approx F_\theta\!\left(r+\frac{\Delta r}{2},\,\theta,\,z\right)
    \left(r+\frac{\Delta r}{2}\right)\Delta\theta,
\]
and along segment~4
\[
  \int_{C_4}\bfF\cdot\hat{\mathbf{t}}\ds
  \approx -F_\theta\!\left(r-\frac{\Delta r}{2},\,\theta,\,z\right)
    \left(r-\frac{\Delta r}{2}\right)\Delta\theta.
\]

Thus,
\begin{align*}
  \frac{1}{\Delta S}\int_{C_2+C_4}\bfF\cdot\hat{\mathbf{t}}\ds
  &\approx \frac{\Delta\theta}{r\,\Delta r\,\Delta\theta}
    \left[\left(r+\frac{\Delta r}{2}\right)
    F_\theta\!\left(r+\frac{\Delta r}{2},\,\theta,\,z\right)\right.\\
  &\qquad\qquad \left.-\left(r-\frac{\Delta r}{2}\right)
    F_\theta\!\left(r-\frac{\Delta r}{2},\,\theta,\,z\right)\right].
\end{align*}

In the limit this becomes $(1/r)(\partial/\partial r)(rF_\theta)$
evaluated at $(r,\theta,z)$.  Hence,
\[
  (\nabla\times\bfF)_z
  \equiv \lim_{\Delta S\to 0}\frac{1}{\Delta S}
  \oint \bfF\cdot\hat{\mathbf{t}}\ds
  = \frac{1}{r}\pd{}{r}(rF_\theta) - \frac{1}{r}\pd{F_r}{\theta}.
\]

Paths for finding the $r$- and $\theta$-components of
$\nabla\times\bfF$ are shown in Figures~III-15(c) and~(d),
respectively.  You are asked to obtain these two components yourself in
Problem~III-8.  For completeness we give all three components of
$\nabla\times\bfF$ in cylindrical coordinates:

\begin{figure}[ht]
\centering
\includegraphics[width=0.6\textwidth]{figures/ch03/fig_III_15c.png}
\caption{}
\label{fig:III-15c}
\end{figure}

\begin{figure}[ht]
\centering
\includegraphics[width=0.6\textwidth]{figures/ch03/fig_III_15d.png}
\caption{}
\label{fig:III-15d}
\end{figure}

\begin{align}
  (\nabla\times\bfF)_r     &= \frac{1}{r}\pd{F_z}{\theta}
                              - \pd{F_\theta}{z}, \notag\\[4pt]
  (\nabla\times\bfF)_\theta &= \pd{F_r}{z} - \pd{F_z}{r},
    \label{eq:cyl-curl}\\[4pt]
  (\nabla\times\bfF)_z     &= \frac{1}{r}\pd{}{r}(rF_\theta)
                              - \frac{1}{r}\pd{F_r}{\theta}. \notag
\end{align}

As an example of calculating the curl in cylindrical coordinates
consider the function
\[
  \bfF(r,\theta,z)
  = \hat{\mathbf{e}}_r\,r^2 z
  + \hat{\mathbf{e}}_\theta\,rz^2\cos\theta
  + \hat{\mathbf{e}}_z\,r^3.
\]
Then
\begin{align*}
  (\nabla\times\bfF)_r
    &= \frac{1}{r}\pd{}{\theta}(r^3) - \pd{}{z}(rz^2\cos\theta)
     = -2rz\cos\theta, \\
  (\nabla\times\bfF)_\theta
    &= \pd{}{z}(r^2 z) - \pd{}{r}(r^3)
     = -2r^2, \\
  (\nabla\times\bfF)_z
    &= \frac{1}{r}\pd{}{r}(rz^2\cos\theta\cdot r)
     - \frac{1}{r}\pd{}{\theta}(r^2 z)
     = 2z^2\cos\theta.
\end{align*}
Hence
\[
  \nabla\times\bfF
  = -2\hat{\mathbf{e}}_r\,rz\cos\theta
    -2\hat{\mathbf{e}}_\theta\,r^2
    +2\hat{\mathbf{e}}_z\,z^2\cos\theta.
\]

The three components of curl~$\bfF$ in \emph{spherical} coordinates
(see Problem~III-9) are as follows:
\begin{align}
  (\nabla\times\bfF)_r
    &= \frac{1}{r\sin\phi}\pd{}{\phi}(\sin\phi\,F_\theta)
     - \frac{1}{r\sin\phi}\pd{F_\phi}{\theta}, \notag\\[4pt]
  (\nabla\times\bfF)_\phi
    &= \frac{1}{r\sin\phi}\pd{F_r}{\theta}
     - \frac{1}{r}\pd{}{r}(rF_\theta), \label{eq:sph-curl}\\[4pt]
  (\nabla\times\bfF)_\theta
    &= \frac{1}{r}\pd{}{r}(rF_\phi)
     - \frac{1}{r}\pd{F_r}{\phi}. \notag
\end{align}

As an example of calculating the curl in spherical coordinates,
consider the function
\[
  \bfF(r,\theta,\phi)
  = \frac{\hat{\mathbf{e}}_r}{r\,\theta}
  + \frac{\hat{\mathbf{e}}_\phi}{r}
  + \frac{\hat{\mathbf{e}}_\theta}{r\cos\phi}.
\]
Then
\begin{align*}
  (\nabla\times\bfF)_r
    &= \frac{1}{r\sin\phi}\pd{}{\phi}
       \!\left(\sin\phi\,\frac{1}{r\cos\phi}\right)
     - \frac{1}{r\sin\phi}\cdot 0
     = \frac{\sec^2\phi}{r^2\sin\phi}, \\[4pt]
  (\nabla\times\bfF)_\phi
    &= \frac{1}{r\sin\phi}\pd{}{\theta}\!\left(\frac{1}{r\,\theta}\right)
     - \frac{1}{r}\pd{}{r}(\cos\phi)
     = -\frac{1}{r^2\theta^2\sin\phi}, \\[4pt]
  (\nabla\times\bfF)_\theta
    &= \frac{1}{r}\pd{}{r}(1)
     - \frac{1}{r}\pd{}{\phi}\!\left(\frac{1}{r\,\theta}\right)
     = 0.
\end{align*}
Hence
\[
  \nabla\times\bfF
  = \frac{\sec^2\phi}{r^2\sin\phi}\,\hat{\mathbf{e}}_r
  - \frac{1}{r^2\theta^2\sin\phi}\,\hat{\mathbf{e}}_\phi.
\]


%% ================================================================
\section{The Meaning of the Curl}
%% ================================================================

The preceding discussion may leave you with the feeling that knowing
how to define and calculate the curl of some vector function is a far
cry from knowing what it is.  The fact that the curl has something to do
with a line integral \emph{around} a closed path (indeed, the word
``curl'' itself) may suggest to you that it somehow has to do with
things rotating, swirling, or curling around.  By means of a few
examples taken from fluid motion, we'll try to make these vague
impressions a little clearer.

\begin{figure}[ht]
\centering
\includegraphics[width=0.5\textwidth]{figures/ch03/fig_III_16.png}
\caption{}
\label{fig:III-16}
\end{figure}

Suppose water is flowing in circular paths, something like the water
draining from a bathtub.  A small volume of the water at a point
$(x,y)$ at time~$t$ has coordinates $x=r\cos\omega t$,
$y=r\sin\omega t$, where $\omega$~is the constant angular velocity of
the water (Figure~III-16).\footnote{This is not a realistic description
of water draining from a tub since rotating water shears tangentially
and its angular velocity will therefore vary with~$r$.  The crude
description we use here is adequate for our purposes and has the virtue
of being simple.}  Thus, its velocity at $(x,y)$ is
\begin{align*}
  \bfv &= \bfi(dx/dt) + \bfj(dy/dt)
        = r\omega[-\bfi\sin\omega t + \bfj\cos\omega t] \\
       &= \omega(-\bfi\,y + \bfj\,x).
\end{align*}

This expression gives what is called the velocity field of the water;
it tells us the velocity of the water at any point $(x,y)$.  Your
intuition probably tells you that, because the motion is circular, this
velocity must have a nonzero curl.  In fact, as you can show very
easily,
\[
  \nabla\times\bfv = 2\bfk\,\omega.
\]

This result should seem quite reasonable because it says that curl of
the velocity is proportional to the angular velocity of the swirling
water.  We see that $\nabla\times\bfv$ is a vector perpendicular to the
plane of motion and in the positive $z$-direction [Figure~III-17(a)].
If the water were rotating around in the other direction, the curl
of~$\bfv$ would then be in the negative $z$-direction
[Figure~III-17(b)].  Note that this is consistent with the right-hand
rule (see page~79--80).  If

\begin{figure}[ht]
\centering
\includegraphics[width=0.4\textwidth]{figures/ch03/fig_III_17a.png}
\caption{}
\label{fig:III-17a}
\end{figure}

\begin{figure}[ht]
\centering
\includegraphics[width=0.4\textwidth]{figures/ch03/fig_III_17b.png}
\caption{}
\label{fig:III-17b}
\end{figure}

\noindent we were to put a small paddle wheel in the water, it would
commence spinning because the impinging water would exert a net torque
on the paddles (Figure~III-18).  Furthermore, the paddle wheel would
rotate with its axis pointing in the direction of the curl.

\begin{figure}[ht]
\centering
\includegraphics[width=0.5\textwidth]{figures/ch03/fig_III_18.png}
\caption{}
\label{fig:III-18}
\end{figure}

Now consider a different velocity field, namely,
\[
  \bfv = \bfj\,v_0\,e^{-y^2/\lambda^2},
\]

\begin{figure}[ht]
\centering
\includegraphics[width=0.5\textwidth]{figures/ch03/fig_III_19.png}
\caption{}
\label{fig:III-19}
\end{figure}

\noindent where $v_0$ and $\lambda$ are constants.  Water with such a
velocity field would have a flow pattern as indicated in Figure~III-19.
The velocity at all points is in the positive $y$-direction, and its
magnitude (indicated by the length of the arrows) varies with~$y$.
Since you see only straight line flow here without any rotational
motion, you would probably guess that $\nabla\times\bfv=0$ in this
case, and you would be right, as a simple calculation shows.  There
would be no net torque on a paddle wheel placed anywhere in this flow
pattern, and as a consequence, it would not
spin.\footnote{If $\nabla\times\bfv=0$, the flow is said to be
irrotational.  Compare with footnote~4.}

Our last example is trickier than the two given previously and shows
that intuition can lead you astray if you're not careful.  Let a
velocity field be given by
\[
  \bfv = \bfj\,v_0\,e^{-x^2/\lambda^2}.
\]

As in the previous example, the velocity in this case is everywhere in
the $y$-direction, but now it varies with~$x$, not~$y$ (Figure~III-20).
Here, as in the preceding example, you see no evidence of rotational
motion and you might guess that $\nabla\times\bfv=0$ once again.  But
as you should show for yourself,
\[
  \nabla\times\bfv
  = -\bfk\,v_0\,\frac{2x}{\lambda^2}\,e^{-x^2/\lambda^2}.
\]

\begin{figure}[ht]
\centering
\includegraphics[width=0.5\textwidth]{figures/ch03/fig_III_20.png}
\caption{}
\label{fig:III-20}
\end{figure}

A small paddle wheel placed in this flow pattern would spin, even
though the water is everywhere moving in the same direction.  The

\begin{figure}[ht]
\centering
\includegraphics[width=0.5\textwidth]{figures/ch03/fig_III_21.png}
\caption{}
\label{fig:III-21}
\end{figure}

\noindent reason this happens is that the velocity of the water varies
with~$x$, so that it strikes one of the paddles ($P$ in
Figure~III-21) with greater velocity than the other~($P'$).  Thus,
there will be a net torque.  In more mathematical terms, the line
integral of $\bfv\cdot\hat{\mathbf{t}}$ around a small rectangle
(Figure~III-22) will be different from zero, for while
\[
  \int_{\text{bottom}}\bfv\cdot\hat{\mathbf{t}}\ds
  = \int_{\text{top}}\bfv\cdot\hat{\mathbf{t}}\ds = 0,
\]
the contributions from the other two sides are
\[
  \int_{\text{right}}\bfv\cdot\hat{\mathbf{t}}\ds
  \approx v_y(x+\Delta x)\,\Delta y
\]

\begin{figure}[ht]
\centering
\includegraphics[width=0.5\textwidth]{figures/ch03/fig_III_22.png}
\caption{}
\label{fig:III-22}
\end{figure}

\noindent and
\[
  \int_{\text{left}}\bfv\cdot\hat{\mathbf{t}}\ds
  \approx -v_y(x)\,\Delta y.
\]

These do not cancel because $v_y(x)\ne v_y(x+\Delta x)$.
Incidentally, you should try to explain to your own satisfaction why in
this example $\nabla\times\bfv$ is in the negative (positive)
$z$-direction when $x$~is positive (negative) and why
$\nabla\times\bfv=0$ at $x=0$.


%% ================================================================
\section{Differential Form of the Circulation Law}
%% ================================================================

The curl is defined to be the limit of circulation to area.  Thus,
\[
  \nhat\cdot\nabla\times\bfE
  = \lim_{\Delta S\to 0}\frac{1}{\Delta S}
    \oint_C \bfE\cdot\hat{\mathbf{t}}\ds,
\]
where $\nhat$~is a unit vector normal to the surface enclosed by~$C$ at
the point about which the curve shrinks to zero.  But if $\bfE$~is an
electrostatic field, then
\[
  \oint_C \bfE\cdot\hat{\mathbf{t}}\ds = 0
\]
for any path~$C$.  It follows that
\[
  \nhat\cdot\nabla\times\bfE = 0.
\]

Since the curve~$C$ is arbitrary, we can arrange matters so that
$\nhat$~is a unit vector pointing in any direction we choose.  Thus,
\begin{alignat*}{3}
  &\text{taking } \nhat=\bfi, &\quad &\text{we have }
    (\nabla\times\bfE)_x = 0; \\
  &\text{taking } \nhat=\bfj, &\quad &\text{we have }
    (\nabla\times\bfE)_y = 0; \\
  &\text{taking } \nhat=\bfk, &\quad &\text{we have }
    (\nabla\times\bfE)_z = 0.
\end{alignat*}

Thus, all three Cartesian components of $\nabla\times\bfE$ vanish and
we can conclude that for an electrostatic field,
\begin{equation}
  \nabla\times\bfE = 0.
\end{equation}

This is the long-sought-after differential form of the circulation law.
We are now in a position to give an alternative and much more tractable
answer to the question ``Can a given vector function $\bfF(x,y,z)$ be
an electrostatic field?''  The answer is
\begin{align*}
  &\text{If } \nabla\times\bfF = 0,
    \text{ then $\bfF$ can be an electrostatic field, and} \\
  &\text{if } \nabla\times\bfF \ne 0,
    \text{ then $\bfF$ cannot be an electrostatic field.}
\end{align*}

This is clearly a much more convenient criterion to apply than our
earlier one (page~77), which required us to determine the line integral
of~$\bfF$ over \emph{all} closed paths!  To see how it works, let us do
several examples.

\emph{Example~1.}  Could $\bfF = K(\bfi\,y + \bfj\,x)$ be an
electrostatic field?  ($K$~is a constant.)  Here we have
\begin{align*}
  \frac{1}{K}\,\nabla\times\bfF
  &= \bfi\!\left(\pd{}{y}\,0-\pd{x}{z}\right)
   + \bfj\!\left(\pd{y}{z}-\pd{}{x}\,0\right)
   + \bfk\!\left(\pd{x}{x}-\pd{y}{y}\right) \\
  &= 0 \;\Rightarrow\; \nabla\times\bfF = 0.
\end{align*}
\emph{Answer:} Yes.

\emph{Example~2.}  Could $\bfF = K(\bfi\,y - \bfj\,x)$ be an
electrostatic field?  In this case
\[
  \frac{1}{K}\,\nabla\times\bfF
  = \bfk\!\left(\pd{(-x)}{x}-\pd{y}{y}\right) = -2\bfk
  \;\Rightarrow\; \nabla\times\bfF = -2\bfk\,K.
\]
\emph{Answer:} No.

From these examples we can see how easy this criterion is to apply.


%% ================================================================
\section{Stokes' Theorem}
%% ================================================================

For the remainder of this chapter we digress from our presentation to
discuss another famous theorem, one strongly reminiscent of the
divergence theorem and yet, as we'll see, quite different from it.
This theorem, named for the mathematician Stokes, relates a line
integral around a closed path to a surface integral over what is called
a \emph{capping surface} of the path, so the first item on our agenda
is to define this term.  Suppose we have a closed curve~$C$, as shown
in Figure~III-23(a), and imagine that it is made of

\begin{figure}[ht]
\centering
\includegraphics[width=0.4\textwidth]{figures/ch03/fig_III_23a.png}
\caption{}
\label{fig:III-23a}
\end{figure}

\noindent wire.  Now let us suppose we attach an elastic membrane to
the wire as indicated in Figure~III-23(b).  This membrane is a capping

\begin{figure}[ht]
\centering
\includegraphics[width=0.4\textwidth]{figures/ch03/fig_III_23b.png}
\caption{}
\label{fig:III-23b}
\end{figure}

\noindent surface of the curve~$C$.  Any other surface which can be
formed by stretching the membrane is also a capping surface; an example
is shown in Figure~III-23(c).  Figure~III-24 shows four different
capping surfaces of a plane circular path: (a) the region of the

\begin{figure}[ht]
\centering
\includegraphics[width=0.4\textwidth]{figures/ch03/fig_III_23c.png}
\caption{}
\label{fig:III-23c}
\end{figure}

\begin{figure}[ht]
\centering
\includegraphics[width=0.7\textwidth]{figures/ch03/fig_III_24.png}
\caption{}
\label{fig:III-24}
\end{figure}

\noindent plane enclosed by the circle, (b)~a hemisphere with the
circle as its rim, (c)~the curved surface of a dunce cap (a right
circular cone), and (d)~the upper and lateral surfaces of a tuna fish
can.

With these preliminary remarks in mind, you won't be surprised to see
us begin this discussion of Stokes' theorem by considering some closed
curve~$C$ and a capping surface~$S$ [Figure~III-25(a)].  As we have
done before, we approximate this capping

\begin{figure}[ht]
\centering
\includegraphics[width=0.5\textwidth]{figures/ch03/fig_III_25a.png}
\caption{}
\label{fig:III-25a}
\end{figure}

\noindent surface by a polyhedron of $N$~faces, each of which is
tangent to $S$ at some point [Figure~III-25(b)].  Note that this will
automatically create a polygon [marked $P$ in Figure~III-25(b)] that is
an approx-

\begin{figure}[ht]
\centering
\includegraphics[width=0.5\textwidth]{figures/ch03/fig_III_25b.png}
\caption{}
\label{fig:III-25b}
\end{figure}

\noindent imation to the curve~$C$.  Let $\bfF(x,y,z)$ be a
well-behaved vector function defined throughout the region of space
occupied by the curve~$C$ and its capping surface~$S$.  Let us form the
circulation of~$\bfF$ around $C_l$, the boundary of the $l$th face of
the polyhedron:
\[
  \oint_{C_l}\bfF\cdot\hat{\mathbf{t}}\ds.
\]

If we do this for each of the faces of the polyhedron and then add
together all the circulations, we assert that this sum will be equal to
the circulation of~$\bfF$ around the polygon~$P$:
\begin{equation}
  \sum_{l=1}^{N}\oint_{C_l}\bfF\cdot\hat{\mathbf{t}}\ds
  = \oint_P\bfF\cdot\hat{\mathbf{t}}\ds.
  \label{eq:III-9}
\end{equation}

This is easy to prove.  Consider two adjacent faces as shown in
Figure~III-26.  The circulation about the face on the left
[Figure~III-26(a)] includes a term from the segment~$AB$, which is
$\int_A^B\bfF\cdot\hat{\mathbf{t}}\ds$.  But the segment~$AB$ is
common to both faces, and its contribution to the circulation around
the \emph{right}-hand face [Figure~III-26(b)] is
\[
  \int_B^A \bfF\cdot\hat{\mathbf{t}}\ds
  = -\int_A^B \bfF\cdot\hat{\mathbf{t}}\ds.
\]

\begin{figure}[ht]
\centering
\includegraphics[width=0.45\textwidth]{figures/ch03/fig_III_26a.png}
\caption{}
\label{fig:III-26a}
\end{figure}

\begin{figure}[ht]
\centering
\includegraphics[width=0.45\textwidth]{figures/ch03/fig_III_26b.png}
\caption{}
\label{fig:III-26b}
\end{figure}

We see that we traverse the common segment~$AB$ one way as part of
the boundary of the left-hand face, and the other way as part of the
boundary of the right-hand face.  Thus, when we add the circulations
of~$\bfF$ over the two faces, the segment~$AB$ contributes
\[
  \int_A^B\bfF\cdot\hat{\mathbf{t}}\ds
  + \int_B^A\bfF\cdot\hat{\mathbf{t}}\ds = 0.
\]

It is clear that any segment common to two adjacent faces contributes
nothing to the sum in Equation~(\ref{eq:III-9}) because such segments
always give rise to pairs of canceling terms.  But \emph{all} segments
are common to pairs of adjacent faces except those that, taken
together, constitute the polygon~$P$.  This establishes
Equation~(\ref{eq:III-9}).

Now we go through an analysis very similar to that which yielded the
divergence theorem.  We write
\begin{equation}
  \oint_P\bfF\cdot\hat{\mathbf{t}}\ds
  = \sum_{l=1}^{N}\oint_{C_l}\bfF\cdot\hat{\mathbf{t}}\ds
  = \sum_{l=1}^{N}\left[\frac{1}{\Delta S_l}
    \oint_{C_l}\bfF\cdot\hat{\mathbf{t}}\ds\right]\Delta S_l,
  \label{eq:III-10}
\end{equation}
where $\Delta S_l$ is the area of the $l$th face.  The quantity in the
square brackets is, approximately, equal to
$\nhat_l\cdot(\nabla\times\bfF)_l$ where $\nhat_l$~is the unit
positive normal on the $l$th face and $(\nabla\times\bfF)_l$~is the
curl of the vector function~$\bfF$ evaluated at the point on the $l$th
face at which it is tangent to~$S$.  We say ``approximately'' because it
is actually the \emph{limit} as $\Delta S_l$ tends to zero of the
bracketed quantity in Equation~(\ref{eq:III-10}), which is to be
identified as $\nhat_l\cdot(\nabla\times\bfF)_l$.  Ignoring this lack
of rigor, we write
\begin{align}
  \lim_{\substack{N\to\infty\\ \text{each }\Delta S_l\to 0}}
  \sum_{l=1}^{N}\left[\frac{1}{\Delta S_l}
    \oint_{C_l}\bfF\cdot\hat{\mathbf{t}}\ds\right]\Delta S_l
  &= \lim_{\substack{N\to\infty\\ \text{each }\Delta S_l\to 0}}
     \sum_{l=1}^{N}\nhat_l\cdot(\nabla\times\bfF)_l\,\Delta S_l
     \notag\\
  &= \iint_S \nhat\cdot\nabla\times\bfF\dS.
  \label{eq:III-11}
\end{align}

Since the curve~$C$ is the limiting shape of the polygon~$P$, we also
have
\begin{equation}
  \lim_{\substack{N\to\infty\\ \text{each }\Delta S_l\to 0}}
  \oint_P\bfF\cdot\hat{\mathbf{t}}\ds
  = \oint_C\bfF\cdot\hat{\mathbf{t}}\ds.
  \label{eq:III-12}
\end{equation}

Combining Equations~(\ref{eq:III-10}), (\ref{eq:III-11}),
and~(\ref{eq:III-12}), we arrive, finally, at Stokes' theorem:
\begin{equation}
  \oint_C\bfF\cdot\hat{\mathbf{t}}\ds
  = \iint_S \nhat\cdot\nabla\times\bfF\dS,
  \label{eq:III-13}
\end{equation}
where $S$~is \emph{any} surface capping the curve~$C$.  Thus, in words,
Stokes' theorem says that the line integral of the tangential component
of a vector function over some closed path equals the surface integral
of the normal component of the curl of that function integrated over
any capping surface of the path.  Stokes' theorem holds for any vector
function~$\bfF$ that is continuous and differentiable and has continuous
derivatives on $C$ and~$S$.

Let's work an example.  Take
$\bfF(x,y,z) = \bfi\,z + \bfj\,x - \bfk\,x$, with $C$~the circle of
radius~1 centered at the origin and lying in the $xy$-plane, and $S$
the part of the $xy$-plane enclosed by the circle [see
Figure~III-27(a)].  Now
\[
  \bfF\cdot\hat{\mathbf{t}}\ds = z\dx + x\dy - x\dz.
\]

\begin{figure}[ht]
\centering
\includegraphics[width=0.5\textwidth]{figures/ch03/fig_III_27a.png}
\caption{}
\label{fig:III-27a}
\end{figure}

Thus, $\oint_C\bfF\cdot\hat{\mathbf{t}}\ds = \oint x\dy$.  Heretofore
we have always parametrized curves with the arc length~$s$.  In this
situation, however, the path~$C$ is most easily parametrized in terms
of the angle~$\theta$ shown in Figure~III-27(b).  Thus, we write
\begin{equation}
  \oint x\dy = \oint x\,\frac{dy}{d\theta}\,d\theta
  = \int_0^{2\pi}\cos^2\theta\,d\theta = \pi,
  \label{eq:III-14}
\end{equation}
where we use $x=\cos\theta$ and $y=\sin\theta$.

\begin{figure}[ht]
\centering
\includegraphics[width=0.5\textwidth]{figures/ch03/fig_III_27b.png}
\caption{}
\label{fig:III-27b}
\end{figure}

Next we calculate
\[
  \nabla\times\bfF
  = \begin{vmatrix}
      \bfi & \bfj & \bfk \\
      \partial/\partial x & \partial/\partial y & \partial/\partial z \\
      z & x & -x
    \end{vmatrix}
  = 2\bfj + \bfk.
\]

The capping surface here is a portion of the $xy$-plane, so that the
unit normal in the positive direction is $\nhat=\bfk$.  Thus,
\[
  \nhat\cdot\nabla\times\bfF = \bfk\cdot(2\bfj+\bfk) = 1
\]
and
\begin{equation}
  \iint_S \nhat\cdot\nabla\times\bfF\dS
  = \iint_S dS = \pi,
  \label{eq:III-15}
\end{equation}
where this last equality follows from the fact that the surface integral
in this case is merely the area of the unit circle.  Since this result
[Equation~(\ref{eq:III-15})] is identical with the one obtained above
[Equation~(\ref{eq:III-14})], we have illustrated Stokes' theorem
[Equation~(\ref{eq:III-13})].

Let's redo this calculation, this time choosing the hemisphere shown
in Figure~III-27(c) as our capping surface~$S$.  Using Equation~(II-13)
with $\bfF$ replaced by $\nabla\times\bfF$, we get

\begin{figure}[ht]
\centering
\includegraphics[width=0.5\textwidth]{figures/ch03/fig_III_27c.png}
\caption{}
\label{fig:III-27c}
\end{figure}

\begin{align*}
  \iint_S \nhat\cdot\nabla\times\bfF\dS
  &= \iint_R \left[-2\!\left(-\frac{y}{z}\right)+1\right]dx\dy \\
  &= 2\iint_R \frac{y}{z}\dx\dy + \iint_R dx\dy
\end{align*}
where $R$~is the unit circle in the $xy$-plane shown in
Figure~III-27(a).  The second integral in the right-most equality above
is just the area of that circle, and so its value is~$\pi$.  The first
integral can be handled by introducing polar coordinates.  We find
\[
  2\iint_R \frac{y}{z}\dx\dy
  = 2\iint_R \frac{y\dx\dy}{\sqrt{1-x^2-y^2}}
  = 2\int_0^{2\pi}\int_0^1
    \frac{r\sin\theta\;r\,dr\,d\theta}{\sqrt{1-r^2}}
  = 2\int_0^{2\pi}\sin\theta\,d\theta
    \int_0^1\frac{r^2\,dr}{\sqrt{1-r^2}}.
\]

It's easy to show that the integral over~$\theta$ is zero.  Hence
$\iint_S\nhat\cdot\nabla\times\bfF\dS=\pi$, consistent with our
earlier result.


%% ================================================================
\section{An Application of Stokes' Theorem}
%% ================================================================

An important application of Stokes' theorem is provided by Amp\`ere's
circuital law.  Consider any closed loop~$C$ enclosing a current~$I$ as
in Figure~III-28.  Note that the direction of~$C$ and that of~$I$
correspond to the same right-hand rule that relates the directions
of~$C$ and the positive normal to a surface capping~$C$.  Amp\`ere's

\begin{figure}[ht]
\centering
\includegraphics[width=0.45\textwidth]{figures/ch03/fig_III_28.png}
\caption{}
\label{fig:III-28}
\end{figure}

\noindent circuital law says that the line integral of the
\emph{magnetic} field~$\bfB$ is related to the current thus:
\[
  \oint_C\bfB\cdot\hat{\mathbf{t}}\ds = \mu_0 I
\]
where the constant $\mu_0$, called the permeability of free space, has
the value $1.257\times 10^{-6}$~newtons per ampere$^2$.  This law, like
Gauss' law and the circulation law, says something about the integral
of a field (the magnetic field in this case), and just as in the two
previous cases, it is convenient to re-express it so that it will tell
us something about the field at a point.  To this end, we first
introduce the current density~$\bfJ$ (see page~52).  Thus, if current
is flowing through an area~$\Delta S$ with normal~$\nhat$
(Figure~III-29), the current density~$\bfJ$ is such that
\[
  \Delta I = \bfJ\cdot\nhat\,\Delta S,
\]

\begin{figure}[ht]
\centering
\includegraphics[width=0.45\textwidth]{figures/ch03/fig_III_29.png}
\caption{}
\label{fig:III-29}
\end{figure}

\noindent where $\Delta I$ is the total current.  That is, current
density is a vector function whose magnitude is the current per unit
area and whose direction is that of the current flow.  If
$\bfJ(x,y,z)$ is the current density, then the total current flowing
through a surface~$S$ is
\[
  \iint_S \bfJ\cdot\nhat\dS.
\]

Thus, Amp\`ere's law can be written
\[
  \oint_C\bfB\cdot\hat{\mathbf{t}}\ds
  = \mu_0\iint_S \bfJ\cdot\nhat\dS.
\]

$S$~can be any surface capping the curve~$C$.  If, as is usually the
case, the current flows through a wire the cross section of which

\begin{figure}[ht]
\centering
\includegraphics[width=0.5\textwidth]{figures/ch03/fig_III_30.png}
\caption{}
\label{fig:III-30}
\end{figure}

\noindent does not include the entire capping surface, it does not
matter; we can integrate over more than the wire cross section if we
remember that $\bfJ\ne 0$ for that part of the surface~$S$ cut by the
wire and $\bfJ=0$ for the rest (Figure~III-30).  Thus,
\[
  \iint_{\substack{\text{cross section}\\\text{of wire}}}
  \bfJ\cdot\nhat\dS
  = \iint_{\substack{\text{entire capping}\\\text{surface }S}}
  \bfJ\cdot\nhat\dS.
\]

Now using Stokes' theorem [Equation~(\ref{eq:III-13})], we have
\[
  \oint_C\bfB\cdot\hat{\mathbf{t}}\ds
  = \iint_S\nhat\cdot\nabla\times\bfB\dS
  = \mu_0\iint_S\nhat\cdot\bfJ\dS.
\]

Since $C$ and $S$ are arbitrary, we conclude that
\[
  \nabla\times\bfB = \mu_0\bfJ.
\]

This is the differential form of Amp\`ere's law.  It is also a special
case of one of Maxwell's equations, valid when the fields do not vary
with time.


%% ================================================================
\section{Stokes' Theorem and Simply Connected Regions}
%% ================================================================

For many purposes, including some important applications, we must be
able to assert that Stokes' theorem holds throughout some region~$D$ in
three-dimensional space.  By this we mean that we want the theorem to
hold for any closed curve~$C$ lying entirely in~$D$ and any capping
surface of~$C$ also lying entirely in~$D$.  This, of course, means the
function~$\bfF$ must be continuous and differentiable and have
continuous first derivatives in~$D$.  But in addition we must impose a
restriction on the region~$D$ itself.  To understand how this comes
about, suppose first that $D$~is the interior of a sphere.  If $\bfF$
is smooth\footnote{Hereafter when we say that a function is ``smooth,''
we'll mean that it is continuous, differentiable, and has continuous
first derivatives.} everywhere in~$D$, then Stokes' theorem holds for
any closed curve~$C$ lying entirely in~$D$, and any capping surface
of~$C$ also lying entirely in~$D$.  In other words, Stokes' theorem
holds everywhere in~$D$.  A little thought should convince you that the
same line of reasoning applies to the region between two concentric
spheres, provided $\bfF$~is smooth in that region.  But for certain
kinds of regions, troubles can arise.  As an example, suppose $D$~is
the interior of a torus (roughly like a bagel or an inflated inner
tube; see Figure~III-31).  The problem in this case is that it's

\begin{figure}[ht]
\centering
\includegraphics[width=0.5\textwidth]{figures/ch03/fig_III_31.png}
\caption{}
\label{fig:III-31}
\end{figure}

\noindent possible to construct a closed curve in~$D$ like the one
shown in the figure with the property that none of its capping surfaces
lies entirely in~$D$.  Although we insist that $\bfF$ be smooth in~$D$,
no conditions are imposed upon it elsewhere, so that outside the region
it may not fulfill the requirements of smoothness that ensure the
validity of Stokes' theorem.  The relation between the line integral
over~$C$ and the surface integral over~$S$ asserted by the theorem can,
and in many cases does, break down if $\bfF$~is not smooth on~$S$.

Mathematicians refer to regions such as the interior of a sphere or
the space between two concentric spheres as \emph{simply connected},
whereas the interior of a torus is not simply connected.  By definition,
a region~$D$ is simply connected if any closed curve lying entirely
in~$D$ can shrink down to a point without leaving~$D$.  Using this
definition, you should be able to verify that the interior of a sphere
and the region between two concentric spheres are both simply connected,
but that the interior of a torus is not.  With the concept of simple
connectedness available to us, we can easily specify the conditions
under which Stokes' theorem holds throughout a region: The vector
function~$\bfF$ must be smooth everywhere in a \emph{simply connected}
region~$D$.  Then Stokes' theorem [Equation~(\ref{eq:III-13})] is valid
for any closed curve~$C$ and any capping surface~$S$ of~$C$, both of
which lie entirely in~$D$.

Most of the time we'll assume that the functions we work with are
smooth and that the regions of interest are simply connected.  There are
situations, however, like the one discussed in the next section, where
simple connectedness plays an essential role, and we'll point them out
as we come to them.


%% ================================================================
\section{Path Independence and the Curl}
%% ================================================================

In our discussion of the differential form of the circulation law, we
showed that because the line integral of an electrostatic field~$\bfE$
is zero over any closed path, the curl of~$\bfE$ is zero.  The same is
true of any vector function~$\bfF$; that is, if
\[
  \oint_C\bfF\cdot\hat{\mathbf{t}}\ds = 0
\]
for all closed paths~$C$, then
\[
  \nabla\times\bfF = 0.
\]

The proof of this fact is precisely the same as the one given on
pages~90--91 with $\bfE$ replaced everywhere by~$\bfF$.

Is the converse of this statement also true?  That is, if
$\nabla\times\bfF=0$, does this imply that the circulation of~$\bfF$ is
zero over all closed paths?  At first glance it might appear that the
answer to this question is yes.  All we have to do is use Stokes'
theorem and observe that since by assumption
$\nabla\times\bfF=0$,
\[
  \oint_C\bfF\cdot\hat{\mathbf{t}}\ds
  = \iint_S \nhat\cdot\nabla\times\bfF\dS = 0.
\]

However, there is a flaw in this line of reasoning.  Recall that the
validity of Stokes' theorem requires that $\bfF$ be smooth \emph{in a
simply connected region}.  If the region is not simply connected,
Stokes' theorem may not hold, at least for some closed paths lying in
the region, and the fact that $\nabla\times\bfF=0$ does not guarantee
that the circulation of~$\bfF$ is zero over all closed paths.  The
closest we can come to a converse is to say that if
$\nabla\times\bfF=0$ \emph{everywhere in a simply connected region},
then the circulation of~$\bfF$ is zero for all closed paths in that
region.  The two statements ``circulation equals zero'' and ``curl
equals zero'' are equivalent only in a simply connected region.

There is a slightly different, but often useful, way to state this
connection between circulation and curl; namely, if
$\int_C\bfF\cdot\hat{\mathbf{t}}\ds$ is independent of path, then
$\nabla\times\bfF=0$, and if $\nabla\times\bfF=0$ in a simply
connected region, then $\int_C\bfF\cdot\hat{\mathbf{t}}\ds$ is
independent of path.  You should have no difficulty in establishing this
for yourself.


%% ================================================================
\begin{problems}
%% ================================================================

\prob
Use an argument like the one given in the text for the Coulomb force
(pages~71--73) to show that $\int_C\bfF\cdot\hat{\mathbf{t}}\ds$ is
independent of path for \emph{any} central force~$\bfF$.

\prob
In the text we obtained the result
\[
  (\nabla\times\bfF)_z = \pd{F_y}{x} - \pd{F_x}{y}
\]
by integrating over a small rectangular path.  As an example of the fact
that this result is independent of the path, rederive it, using the
triangular path shown in the figure.

\begin{figure}[ht]
\centering
\includegraphics[width=0.4\textwidth]{figures/ch03/fig_prob_III_2.png}
\caption{}
\label{fig:prob-III-2}
\end{figure}

\prob
Calculate the curl of each of the following functions using
Equation~(\ref{eq:III-7b}):
\begin{enumerate}[label=(\alph*)]
\item $\bfi\,z^2 + \bfj\,x^2 - \bfk\,y^2$.
\item $3\bfi\,xz - \bfk\,x^2$.
\item $\bfi\,e^{-y} + \bfj\,e^{-z} + \bfk\,e^{-x}$.
\item $\bfi\,yz + \bfj\,xz + \bfk\,xy$.
\item $-\bfi\,yz + \bfj\,xz$.
\item $\bfi\,x + \bfj\,y + \bfk(x^2+y^2)$.
\item $\bfi\,xy + \bfj\,y^2 + \bfk\,yz$.
\item $(\bfi\,x+\bfj\,y+\bfk\,z)/(x^2+y^2+z^2)^{3/2}$,
      $(x,y,z)\ne(0,0,0)$.
\end{enumerate}

\prob
\begin{enumerate}[label=(\alph*)]
\item Calculate $\oint\bfF\cdot\hat{\mathbf{t}}\ds$ for the function
in Problem~III-3(a) over a square path of side~$s$ centered at
$(x_0,y_0,0)$, lying in the $xy$-plane, and oriented so that each side
is parallel to the $x$- or $y$-axis.
\item Divide the result of part~(a) by the area of the square and take
the limit of the quotient as $s\to 0$.  Compare your result with the
$z$-component of the curl found in Problem~III-3(a).
\item Repeat parts~(a) and~(b) for the functions in Problem~III-3(b),
(c), and~(d).  (You may find it interesting to try paths of different
orientations and/or shapes.)
\end{enumerate}

\prob
\begin{enumerate}[label=(\alph*)]
\item Calculate $\oint\bfF\cdot\hat{\mathbf{t}}\ds$ where
\[
  \bfF = \bfk(y+y^2)
\]
over the perimeter of the triangle shown in the figure (integrate in the
direction indicated by the arrows).

\begin{figure}[ht]
\centering
\includegraphics[width=0.5\textwidth]{figures/ch03/fig_prob_III_5.png}
\caption{}
\label{fig:prob-III-5}
\end{figure}

\item Divide the result of part~(a) by the area of the triangle and
take the limit as $a\to 0$.
\item Show that the result of part~(b) is
$\nhat\cdot\nabla\times\bfF$ evaluated at $(0,0,0)$ where $\nhat$~is
the unit vector normal to the triangle and directed away from the
origin.
\end{enumerate}

\prob
Show that
\[
  \nabla\times\frac{\bfA\times\bfr}{2} = \bfA,
\]
where $\bfr = \bfi\,x + \bfj\,y + \bfk\,z$ and $\bfA$~is a constant
vector.

\prob
Show that $\nabla\cdot(\nabla\times\bfF) = 0$.  (Assume that mixed
second partial derivatives are independent of the order of
differentiation.  For example,
$\partial^2 F_z/\partial x\,\partial z =
\partial^2 F_z/\partial z\,\partial x$.)

\prob
In the text (pages~82--86) we obtained the $z$-component of
$\nabla\times\bfF$ in cylindrical coordinates.  Proceeding in the same
way, obtain the $\theta$- and $r$-components given on page~86.

\prob
Following the procedure suggested in the text (pages~82--86), obtain
the expression for $\nabla\times\bfF$ in spherical coordinates given on
page~86.  The figures given on page~106 will be helpful.

\prob
\begin{enumerate}[label=(\alph*)]
\item Rewrite the function in Problem~III-3(e) in cylindrical
coordinates and compute its curl using the expression given on
page~86.  Convert your result back to Cartesian coordinates and compare
with the answer obtained in Problem~III-3(e) (see Problem~II-16).
\item Repeat the above calculation for the function of
Problem~III-3(f).
\end{enumerate}

\prob
\begin{enumerate}[label=(\alph*)]
\item Rewrite the function in Problem~III-3(g) in spherical coordinates
and compute its curl using the expression given on page~86.  Convert
your result back to Cartesian coordinates and compare with the answer
obtained in Problem~III-3(g) (see Problem~II-17).
\item Repeat the above calculation for the function of
Problem~III-3(h).
\end{enumerate}

\prob
Any central force can be written in the form
\[
  \bfF(r) = \hat{\mathbf{e}}_r\,f(r),
\]
where $\hat{\mathbf{e}}_r$ is a unit vector in the radial direction and
$f$~is a scalar function.  Show by direct calculation of the curl that
this function is irrotational (that is, $\nabla\times\bfF=0$).

\prob
Which of the functions in Problem~III-3 could be electrostatic fields?

\prob
Use Stokes' theorem to show that
\[
  \oint_C \hat{\mathbf{t}}\ds = 0,
\]
where $C$~is a closed curve and $\hat{\mathbf{t}}$~is a unit vector
tangent to the curve~$C$.

\prob
Verify Stokes' theorem
\[
  \oint_C\bfF\cdot\hat{\mathbf{t}}\ds
  = \iint_S \nhat\cdot\nabla\times\bfF\dS
\]
in each of the following cases:
\begin{enumerate}[label=(\alph*)]
\item $\bfF = \bfi\,z^2 - \bfj\,y^2$.

$C$, the square of side~1 lying in the $xz$-plane and directed as
shown.

$S$, the five squares $S_1$, $S_2$, $S_3$, $S_4$, and $S_5$ as shown
in the figure.

\begin{figure}[ht]
\centering
\includegraphics[width=0.5\textwidth]{figures/ch03/fig_prob_III_15a.png}
\caption{}
\label{fig:prob-III-15a}
\end{figure}

\item $\bfF = \bfi\,y + \bfj\,z + \bfk\,x$.

$C$, the three quarter circle arcs $C_1$, $C_2$, and $C_3$ directed as
shown in the figure.

$S$, the octant of the sphere $x^2+y^2+z^2=1$ enclosed by the three
arcs.

\begin{figure}[ht]
\centering
\includegraphics[width=0.5\textwidth]{figures/ch03/fig_prob_III_15b.png}
\caption{}
\label{fig:prob-III-15b}
\end{figure}

\item $\bfF = \bfi\,y - \bfj\,x + \bfk\,z$.

$C$, the circle of radius~$R$ lying in the $xy$-plane, centered at
$(0,0,0)$ and directed as shown in the figure.

$S$, the curved and upper surfaces of the cylinder of radius~$R$ and
height~$h$.

\begin{figure}[ht]
\centering
\includegraphics[width=0.5\textwidth]{figures/ch03/fig_prob_III_15c.png}
\caption{}
\label{fig:prob-III-15c}
\end{figure}
\end{enumerate}

\prob
\begin{enumerate}[label=(\alph*)]
\item Consider a vector function with the property
$\nabla\times\bfF=0$ everywhere on two closed curves $C_1$ and $C_2$
and on any capping surface~$S$ of the region enclosed by them (see the
figure).  Show that the circulation of~$\bfF$ around $C_1$ equals the
circulation of~$\bfF$ around~$C_2$.  In calculating the circulations
direct the curves as indicated by the arrows in the figure.

\begin{figure}[ht]
\centering
\includegraphics[width=0.45\textwidth]{figures/ch03/fig_prob_III_16a.png}
\caption{}
\label{fig:prob-III-16a}
\end{figure}

\item The magnetic field due to an infinitely long straight wire
carrying a uniform current~$I$ is
$\bfB=(\mu_0 I/2\pi r)\hat{\mathbf{e}}_\theta$.  Show that
$\nabla\times\bfB=0$ everywhere except at $r=0$.

\begin{figure}[ht]
\centering
\includegraphics[width=0.45\textwidth]{figures/ch03/fig_prob_III_16b.png}
\caption{}
\label{fig:prob-III-16b}
\end{figure}

\item Prove Amp\`ere's circuital law for the field of the wire given in
part~(b).  [\emph{Hint:} Use the result of~(b) to find the circulation
of~$\bfB$ around a circle with the wire passing through its center and
normal to its plane.  Then use the result of part~(a) to relate this
circulation to the circulation around an arbitrary curve enclosing the
current.]
\end{enumerate}

\prob
\begin{enumerate}[label=(\alph*)]
\item Consider the function given in cylindrical coordinates by
\[
  \bfF(r,\theta,z) = \frac{\hat{\mathbf{e}}_\theta}{r}.
\]
Show that Stokes' theorem does not hold for this function if $C$~is the
circle of radius~$R$ in the $xy$-plane centered at the origin, and
$S$~is the portion of the $xy$-plane enclosed by~$C$.  Why does the
theorem fail in this case?
\item Consider the region~$D$ that consists of all of
three-dimensional space with the $z$-axis removed.  Is the function
$\bfF$ defined in~(a) smooth in~$D$?  Does Stokes' theorem hold
in~$D$?  Is $D$~a simply connected region?
\end{enumerate}

\prob
The electromotive force $\mathscr{E}$ in a circuit~$C$ is equal to the
circulation of the electric field~$\bfE$ around the circuit:
\[
  \mathscr{E} = \oint_C\bfE\cdot\hat{\mathbf{t}}\ds.
\]
Faraday discovered that in a stationary circuit an electromotive force
is induced by a changing magnetic flux.  That is,
\[
  \mathscr{E} = -\frac{d\Phi}{dt},
\]
where
\[
  \Phi = \iint_S \bfB\cdot\nhat\dS,
\]
$t$~is time (don't confuse it with the tangent vector
$\hat{\mathbf{t}}$), and $S$~is any capping surface of~$C$.  Use this
information and Stokes' theorem to derive the equation
\[
  \nabla\times\bfE = -\pd{\bfB}{t},
\]
which is one of Maxwell's equations.

\prob
Determine the value of the line integral
$\int_C\bfF\cdot\hat{\mathbf{t}}\ds$, where
\[
  \bfF = (e^{-y}-ze^{-x})\bfi + (e^{-z}-xe^{-y})\bfj
       + (e^{-x}-ye^{-z})\bfk
\]
and $C$~is the path
\[
  x = \frac{1}{\ln 2}\ln(1+p),\quad
  y = \sin\frac{\pi p}{2},\quad
  z = \frac{1-e^p}{1-e},
  \qquad 0\le p\le 1
\]
from $(0,0,0)$ to $(1,1,1)$.  [\emph{Suggestion:} Think before you
write!]

\prob
Maxwell's equations are
\[
  \nabla\cdot\bfE = \rho/\epszero,\qquad
  \nabla\cdot\bfB = 0,
\]
\[
  \nabla\times\bfE = -\pd{\bfB}{t},\qquad
  \nabla\times\bfB = \epszero\mu_0\pd{\bfE}{t} + \mu_0\bfJ,
\]
where $\bfE$~is the electric field, $\bfB$~the magnetic field,
$\rho$~the charge density, and $\bfJ$~the current density.  Use
Maxwell's equations to derive the continuity equation
\[
  \nabla\cdot\bfJ + \pd{\rho}{t} = 0.
\]
Interpret this equation.

\prob
The electromagnetic field stores energy, and it is possible to show
that in a volume~$V$ the amount of electromagnetic energy is
\[
  \iiint_V \rho_E\dV,
\]
where the energy density
\[
  \rho_E = \tfrac{1}{2}(\epszero\bfE\cdot\bfE
          + \bfB\cdot\bfB/\mu_0)
         = \tfrac{1}{2}(\epszero E^2 + B^2/\mu_0).
\]
Use Maxwell's equations (see Problem~III-20) to show that
\[
  \pd{\rho_E}{t}
  + \nabla\cdot\!\left(\frac{\bfE\times\bfB}{\mu_0}\right)
  = -\bfJ\cdot\bfE.
\]
Interpret this equation.

\prob
\begin{enumerate}[label=(\alph*)]
\item Apply the divergence theorem to the function
\[
  \bfG(x,y) = \bfi\,G_x(x,y) + \bfj\,G_y(x,y),
\]
using for $V$ and $S$ the volume and surface shown in the diagram; its
bottom is a region~$R$ of the $xy$-plane, its top has the same shape as,
and is parallel to, the bottom, and its side is parallel to the
$z$-axis.  In this way obtain the relation
\[
  \oint_C G_x\dy - G_y\dx
  = \iint_R\left(\pd{G_x}{x}+\pd{G_y}{y}\right)dx\dy,
\]
which is the divergence theorem in two dimensions.

\begin{figure}[ht]
\centering
\includegraphics[width=0.5\textwidth]{figures/ch03/fig_prob_III_22.png}
\caption{}
\label{fig:prob-III-22}
\end{figure}

\item Apply Stokes' theorem to the function
\[
  \bfF(x,y) = \bfi\,F_x(x,y) + \bfj\,F_y(x,y)
\]
using for $C$~a closed curve lying entirely in the $xy$-plane and for
$S$~the region~$R$ of the $xy$-plane enclosed by~$C$.  In this way
obtain the relation
\[
  \oint_C F_x\dx + F_y\dy
  = \iint_R\left(\pd{F_y}{x}-\pd{F_x}{y}\right)dx\dy,
\]
which is Stokes' theorem in two dimensions.
\item Show that in two dimensions the divergence theorem and Stokes'
theorem are identical.
\end{enumerate}

\prob
\begin{enumerate}[label=(\alph*)]
\item Let $C$~be a closed curve lying in the $xy$-plane.  What
condition must the function~$\bfF$ satisfy in order that
\[
  \oint_C\bfF\cdot\hat{\mathbf{t}}\ds = A,
\]
where $A$~is the area enclosed by~$C$?  [\emph{Hint:} See
Problem~III-22.]
\item Give some examples of functions~$\bfF$ having the property
described in~(a).
\item Use line integrals to find formulas for the area of
\begin{enumerate}[label=(\roman*)]
  \item a rectangle.
  \item a right triangle.
  \item a circle.
\end{enumerate}
\item Show that the area enclosed by the plane curve~$C$ is the
magnitude of
\[
  \frac{1}{2}\oint_C \bfr\times\hat{\mathbf{t}}\ds,
\]
where $\bfr = \bfi\,x + \bfj\,y$.
\end{enumerate}

\prob
\begin{enumerate}[label=(\alph*)]
\item There is an important theorem in vector calculus that says
$\nabla\cdot\bfG=0$ (where $\bfG$~is some differentiable vector
function) implies and is implied by $\bfG=\nabla\times\mathbf{H}$
(where $\mathbf{H}$~is another differentiable function).  To prove this
we note first of all that $\bfG=\nabla\times\mathbf{H}$ implies that
$\nabla\cdot\bfG=0$ (see Problem~III-7).  To show that
$\nabla\cdot\bfG=0$ implies that we can write
$\bfG=\nabla\times\mathbf{H}$, the simplest procedure is to give
$\mathbf{H}$:
\begin{align*}
  H_x &= 0, \\
  H_y &= \int_{x_0}^{x}G_z(x',y,z)\,dx', \\
  H_z &= -\int_{x_0}^{x}G_y(x',y,z)\,dx'
        + \int_{y_0}^{y}G_x(x_0,y',z)\,dy',
\end{align*}
where $x_0$ and $y_0$ are arbitrary constants.  Show by direct
calculation that if $\nabla\cdot\bfG=0$, then
$\bfG=\nabla\times\mathbf{H}$.
\item Is the vector function~$\mathbf{H}$ specified in~(a) unique?
That is, can we alter it in any way without invalidating the relation
$\bfG=\nabla\times\mathbf{H}$?
\end{enumerate}

\prob
Determine in which of the following cases it is possible to write
$\bfG=\nabla\times\mathbf{H}$.  In the cases where it is possible,
find~$\mathbf{H}$ (see Problem~III-24).
\begin{enumerate}[label=(\alph*)]
\item $\bfG = \bfi\,y + \bfj\,z + \bfk\,x$.
\item $\bfG = B_0\bfk$, \;$B_0$ a constant.
\item $\bfG = \bfi\,x^2 - \bfk\,y^2$.
\item $\bfG = 2\bfi\,x - \bfj\,y - \bfk\,z$.
\item $\bfG = 2\bfi\,x - \bfj\,y + \bfk\,z$.
\end{enumerate}

\prob
Since the divergence of any magnetic field~$\bfB$ is zero, we can
write $\bfB=\nabla\times\bfA$ (see Problem~III-24).  Prove that the
circulation of~$\bfA$ around an arbitrary closed path~$C$ is equal to
the flux of~$\bfB$ through any surface~$S$ capping~$C$.

\prob
Prove the statement made in Problem~III-24(a) by applying Stokes'
theorem and the divergence theorem.  [\emph{Hint:} See the diagram
below.]

\begin{figure}[ht]
\centering
\includegraphics[width=0.45\textwidth]{figures/ch03/fig_prob_III_27.png}
\caption{}
\label{fig:prob-III-27}
\end{figure}

\prob
\begin{enumerate}[label=(\alph*)]
\item What is the integral form of the equation
$\bfG=\nabla\times\mathbf{H}$?  [\emph{Hint:} Compare the differential
and integral forms of Amp\`ere's circuital law.]
\item Verify your result in part~(a) using for $\bfG$ and $\mathbf{H}$
functions selected from Problem~III-25, and paths and surfaces of
integration of your own choice.
\end{enumerate}

\prob
In the text we defined the curl as the limit of a certain ratio.  An
alternative definition is provided by the equation
\[
  \nabla\times\bfF
  = \lim_{\Delta V\to 0}\frac{1}{\Delta V}
    \iint_S \nhat\times\bfF\dS,
\]
where $\bfF$~is a vector function of position, the integration is
carried out over a closed surface~$S$ which encloses the
volume~$\Delta V$, and $\nhat$~is the unit vector normal to $S$
pointing outward from the enclosed volume.  (This definition does not
display the geometric significance of the curl as well as the one given
in the text.  Nonetheless, in one respect at least it may be
preferable: it gives the $\nabla\times\bfF$ rather than just a
component of it.)
\begin{enumerate}[label=(\alph*)]
\item Following a procedure similar to the one used in the text in
treating the divergence, integrate over a ``cuboid'' and show that the
definition given above yields Equation~(\ref{eq:III-7b}).
\item Arguing as we did in the text in establishing the divergence
theorem, use the above expression for the curl to derive the equation
\[
  \iint_S \nhat\times\bfF\dS = \iiint_V\nabla\times\bfF\dV,
\]
where $V$~is the volume enclosed by~$S$.
\item Derive the equation of part~(b) directly from the divergence
theorem.  [\emph{Hint:} In the divergence theorem [Equation~(II-30)]
replace $\bfF$ by $\mathbf{e}\times\bfF$, where $\mathbf{e}$~is an
arbitrary constant vector.]
\item Verify the equation of part~(b) for
$\bfF=\bfi\,y-\bfj\,z+\bfk\,x$ and $V$~the unit cube shown in the
figure.

\begin{figure}[ht]
\centering
\includegraphics[width=0.45\textwidth]{figures/ch03/fig_prob_III_29.png}
\caption{}
\label{fig:prob-III-29}
\end{figure}
\end{enumerate}

\prob
The result
\[
  (\nabla\times\bfF)_z = \pd{F_y}{x} - \pd{F_x}{y}
\]
has been established by calculating the circulation of~$\bfF$ around a
rectangle (see the text, pages~75~ff.) and around a right triangle (see
Problem~III-2).  In this problem you will show that the result holds
when the circulation is calculated around \emph{any} closed curve lying
in the $xy$-plane.
\begin{enumerate}[label=(\alph*)]
\item Approximate an arbitrary closed curve~$C$ in the $xy$-plane by a
polygon~$P$ as shown in the figure.  Subdivide the area enclosed by~$P$
into $N$~patches of which the $l$th has area~$\Delta S_l$.  Convince
yourself by means of a sketch that this subdivision can be made with
only two kinds of patches: rectangles and right triangles.

\begin{figure}[ht]
\centering
\includegraphics[width=0.45\textwidth]{figures/ch03/fig_prob_III_30.png}
\caption{}
\label{fig:prob-III-30}
\end{figure}

\item Letting $C(x,y) = \partial F_y/\partial x - \partial
F_x/\partial y$, use Taylor series to show that for $N$ large and each
$\Delta S_l$ small,
\begin{align*}
  \oint_P\bfF\cdot\hat{\mathbf{t}}\ds
  &= \sum_{l=1}^{N}\oint_{C_l}\bfF\cdot\hat{\mathbf{t}}\ds \\
  &\cong C(x_0,y_0)\,\Delta A
   + \left(\pd{C}{x}\right)_{x_0,y_0}
     \sum_{l=1}^{N}(x_l-x_0)\,\Delta S_l \\
  &\quad + \left(\pd{C}{y}\right)_{x_0,y_0}
     \sum_{l=1}^{N}(y_l-y_0)\,\Delta S_l + \cdots,
\end{align*}
where $C_l$~is the perimeter of the $l$th patch, $(x_0,y_0)$~is some
point in $P$, and $\Delta A$~is the area enclosed by~$P$.
\item Show that
\[
  \lim_{\substack{N\to\infty\\\text{each }\Delta S_l\to 0}}
  \oint_P\bfF\cdot\hat{\mathbf{t}}\ds
  = \oint_C\bfF\cdot\hat{\mathbf{t}}\ds
\]
where $\Delta S$~is the area of the region~$R$ enclosed by~$C$ and
$(\bar{x},\bar{y})$ are the coordinates of the centroid of the
region~$R$; that is,
\[
  \bar{x} = \frac{1}{\Delta S}\iint_R x\dx\dy
  \qquad\text{and}\qquad
  \bar{y} = \frac{1}{\Delta S}\iint_R y\dx\dy.
\]
\item Finally, calculate
\[
  (\nabla\times\bfF)_z
  = \lim_{\substack{\Delta S\to 0 \\ \text{about }x_0,y_0}}
    \frac{1}{\Delta S}\oint_C\bfF\cdot\hat{\mathbf{t}}\ds.
\]
\end{enumerate}

\end{problems}
